\section{Project Management}
\label{sec:methodology_project_management}

\subsection{Development Process}
\label{subsec:methodology_development_process}

The study followed an iterative development process rather than a strictly linear sequence of tasks.
The initial proposal was kept sufficiently broad to allow the research direction, experimental design, and technical emphasis to be refined as preliminary results clarified the behavior of the GLPE exploration method.
Across successive cycles, the method specification, experimental procedure, analysis strategy, and thesis manuscript were revised to maintain consistency between the implemented study and its academic presentation.

Development and experiment execution preceded the final writing of the thesis manuscript.
This ordering allowed the final methodology and discussion to be based on completed experiments rather than on planned procedures alone.
The final manuscript was then written to consolidate the theoretical framing, experimental protocol, results, and interpretation into a standalone thesis.

\subsection{Budget and Cost Management}
\label{subsec:methodology_costs}

Cloud compute was the only direct project cost incurred for the experimental component.
The recorded expenditure is summarized in Table~\ref{tab:project_management_budget}.
The amount aggregates the cloud resources used for preliminary experiments, exploratory tests, development runs, and final experimental runs, and should be interpreted as a project-level compute cost rather than as a per-experiment budget.

\begin{table}[htbp]
\centering
\caption{Cloud compute cost for the study}
\label{tab:project_management_budget}
\begin{tabular}{p{0.32\linewidth}p{0.43\linewidth}r}
\toprule
Cost item & Coverage & Cost \\
\midrule
Cloud compute & Preliminary experiments, exploratory tests, development runs, and final experimental runs & \(\mathrm{PHP}\,41{,}266.87\) \\
\midrule
\textbf{Total} & \textbf{Recorded cloud compute expenditure} & \(\mathbf{\mathrm{PHP}\,41{,}266.87}\) \\
\bottomrule
\end{tabular}
\end{table}

No separate hardware procurement cost was incurred for the experimental component.

\subsection{Timeline}
\label{subsec:methodology_timeline}

The study was conducted from February 2025 to May 2026.
The timeline involved overlapping activities, including proposal preparation, method refinement, implementation, experiment execution, analysis, writing, and presentation preparation.
These activities were not strictly separated because results from one stage often informed revisions to another stage.

The early phase from February 2025 focused on defining the research problem, surveying relevant literature, and preparing a proposal that allowed subsequent refinement of the technical direction.
Main development and final experimental work began around December 2025.
The main experimental runs were completed in December 2025, after which the remaining work focused on organizing the results, completing the analysis, writing the thesis manuscript, and preparing presentation materials.
Some intervals involved reduced activity because of academic and scheduling constraints, but the overall work continued through iterative refinement until completion in May 2026.

\subsection{Responsibilities}
\label{subsec:methodology_responsibilities}

Mikael Vincent Tumampos handled the primary research work, including the formulation of the method, implementation, experiment execution, result analysis, thesis writing, and preparation of the technical presentation.
Ervin Joshua Guirnela supported the completion of the thesis by assisting in the preparation and organization of presentation materials, coordinating academic and institutional requirements, and helping present the study in a form suitable for review and defense.
Christine Peña served as the thesis adviser and provided guidance on institutional requirements, thesis structure, academic presentation, and compliance with undergraduate thesis standards.